\chapter {Pong as a Testbed}
\label{chap:pong}
Pong provides a controlled setting in which to compare the two paradigms developed in the previous two chapters. The environment is complex enough to require learning visual representations from raw pixels, yet simple enough that both \acrshort{dqn} and \acrshort{a2c} can be trained to play on a single GPU reasonably. This chapter formalizes Pong as a Markov decision process, describes the preprocessing pipeline that transforms raw frames into a suitable state representation, and presents the implementation and training results for each method.

\section{Environment Description and Preprocessing Pipeline}
Pong is a game containing two paddles and one ball in a court. The agent controls a paddle and the other one is controlled by a built-in opponent. A round ends when the balls passes one paddle. A player wins when 21 points are scored.

The mapping of Pong to the finite \acrshort{mdp} in \eqref{eq:finite_mdp} is straightforward. The state space $S$ consists of the raw pixel image of 210x160x3 values returned by the Gymnasium \acrshort{ale} wrappers at each timestamp. This embeds the ball position, velocity, paddle positions and the score which are not directly accessible and must be learned by the agent. The action space $A$ is discrete and consist of only 6 values NOOP (does nothing), FIRE (does nothing), RIGHT (go up), LEFT (go down), RIGHTFIRE (go up), LEFTFIRE (go down). The transition dynamics $p$ is based on deterministic physics and is defined in the \acrshort{ale} environment. The agent has no access to $p$ which must be learned through interaction making it a model-free environment. The reward signal $R$ is sparse where the agent receives a +1 when it scores, a -1 when the opponent scores and 0 otherwise. This sparsity makes Pong a meaningful testbed, testing whether an algorithm can propagate reward backward through long sequences of actions. The discount factor $\gamma$ is fixed by the environment at 0.99.

The \acrshort{mdp} formalism assumes that the environment follows the Markov property of encapsulating all that is required to act optimally. Simply put, the observations should allow the states to be distinguishable from one another. In the case of Pong, a single frame is not enough to learn the speed and direction of objects like the ball and the paddle. This type of environment is called a partially observable \acrshort{mdp} \citep{kaelblingPlanningActingPartially1998}. 

\section{DQN Implementation and Results}

\section{Policy Gradient Implementation and Results}
\label{sec:a2c_results}